\originalsection{18.112历史期末考试（2006）}
\sourceinfo{题面finalprob.pdf, PDF pp.2--3实印“Dec. 20, 2006”, 开卷; 官方解答final.pdf, pp.2--5. MIT OCW 18.112 Fall 2008归档, CC BY-NC-SA 4.0. 6题, 不将2008归档日期当作卷面年份.}
\begin{exercise}\label{ass:f112-1}
\textbf{原题 F112.1（20分）。} 复数 $a,b,c$ 满足 $(b-a)/(c-a)=(a-c)/(b-c)$. 证明顶点为这三数的三角形是等边三角形, 即三边模相等.
\end{exercise}
\begin{solution}
据官方解答的代数路线补全. 分式有定义要求 $c\ne a,b\ne c$; 若 $b=a$ 则右侧非零而左侧0, 矛盾, 故三点不同. 令 $t=(b-a)/(c-a)$, 则右侧为 $-1/(t-1)$, 条件给 $t^2-t+1=0$, 所以 $t=(1\pm\ii\sqrt3)/2$, $|t|=|t-1|=1$. 因 $b-a=t(c-a)$, $b-c=(t-1)(c-a)$, 得所求三边相等. 非零虚部同时证明三点不共线.
\end{solution}
\begin{exercise}\label{ass:f112-2}
\textbf{原题 F112.2（15分）。} 判断 $\sum_{n=1}^{\infty}z^n/(1+z^{2n})$ 在何处收敛, 以及其和 $f(z)$ 在何处全纯, 并说明理由.
\end{exercise}
\begin{solution}
据官方解答翻译并明确校正其边界常数、紧集估计方向和0点处理. 当 $|z|<1$, 对任意紧子集取 $r<1$ 包含它, 用逆三角不等式得 $|z^n/(1+z^{2n})|\le r^n/(1-r^{2n})\le r^n/(1-r^2)$. 几何级数给绝对局部一致收敛, 和全纯, 且0处每项0, 故 $f(0)=0$.

当 $|z|>1$, 先将项改写成 $z^{-n}/(1+z^{-2n})$. 在紧子集上取 $R>1$ 为模的下界, 得上界 $R^{-n}/(1-R^{-2n})\le R^{-n}/(1-R^{-2})$, 同样局部一致收敛且和全纯.

当 $|z|=1$, 写 $z=\ee^{\ii\theta}$. 若有某个 $n$ 使分母0, 原级数在该点无定义. 否则项为 $1/(2\cos n\theta)$, 绝对值至少 $1/2$, 不趋0, 故不收敛. 官方在此漏因子2; 其紧集界用了 $C^n+C^{-n}$ 作分母, 对复数不成立, 必须用上述逆三角界. 故收敛与由此级数定义的全纯域恰为单位圆内、外两个开集; 没有从本题额外断言解析延拓不存在.
\end{solution}
\begin{exercise}\label{ass:f112-3}
\textbf{原题 F112.3（15分）。} $\gamma$ 为半径2、中心0的正向圆, 求 $\int_\gamma |z|\ee^z/z^2\dd z$.
\end{exercise}
\begin{solution}据官方解答翻译. 只在积分路径上用 $|z|=2$, 改为 $\int_\gamma2\ee^z/z^2\dd z$. 后者0处留数2, 积分 $4\pi\ii$. 这不意味着原含 $|z|$ 的函数在圆盘内解析; 替换依赖路径恒模.\end{solution}
\begin{exercise}\label{ass:f112-4}
\textbf{原题 F112.4（15分）。} $f$ 在 $z_0$ 有 $h$ 阶极点, 证明
\[
\Res_{z_0}f=\frac1{(h-1)!}\left.\frac{\mathrm d^{h-1}}{\mathrm dz^{h-1}}\big((z-z_0)^hf(z)\big)\right|_{z=z_0}.
\]
\end{exercise}
\begin{solution}据官方解答翻译. $g=(z-z_0)^hf$ 在 $z_0$ 可去延拓为解析函数. Taylor 展开 $g=\sum_{k\ge0}g^{(k)}(z_0)(z-z_0)^k/k!$, 除以 $(z-z_0)^h$ 后 $f$ 的 $-1$ 次系数正是 $k=h-1$ 的系数, 即所求留数. 导数在延拓后的 $g$ 上取, 不能对未定义点直接代入原分式.\end{solution}
\begin{exercise}\label{ass:f112-5}
\textbf{原题 F112.5（20分）。} 用几何级数求 $f(z)=1/((z-1)(z-2))$ 分别在 $|z|<1$ 与 $|z|>2$ 的 Laurent 展开.
\end{exercise}
\begin{solution}
据官方解答翻译. $f=-1/(z-1)+1/(z-2)$. 在内侧分别展开得 $f=\sum_{n\ge0}(1-2^{-n-1})z^n=1/2+3z/4+7z^2/8+\cdots$. 在外侧展开得 $f=\sum_{n\ge0}(2^n-1)z^{-n-1}=z^{-2}+3z^{-3}+7z^{-4}+\cdots$. 第一个表达式由 $(1-z)^{-1}-(2-z)^{-1}$ 推出, 第二个由 $-z^{-1}(1-1/z)^{-1}+z^{-1}(1-2/z)^{-1}$ 推出, 收敛条件正是题定区域.
\end{solution}
\begin{exercise}\label{ass:f112-6}
\textbf{原题 F112.6（15分）。} $f$ 在 $|z|\le1$ 上解析且圆周上 $|f(z)|<1$, 证明 $f(z)-z=0$ 在圆内恰有一个解.
\end{exercise}
\begin{solution}据官方解答翻译. 将“闭圆盘上解析”理解为其开邻域解析. 圆周上 $|f|<|-z|$, Rouché 定理使 $f-z$ 与 $-z$ 有相同内部零点数1（计重数）. 故唯一不动点在内部且为简单零点. 与18.04期末模拟题2(c)完全重复, 不在去重统计中再算独立题意, 但仍保留各卷完整编号.\end{solution}
